\documentclass[11pt]{article}

\usepackage[margin=1in]{geometry}
\usepackage{amsmath,amssymb,amsthm}
\usepackage{array}
\usepackage{graphicx}
\usepackage[colorlinks=true,linkcolor=blue,urlcolor=blue,citecolor=blue]{hyperref}
\emergencystretch=2em

\newtheorem{theorem}{Theorem}
\newtheorem{corollary}[theorem]{Corollary}
\theoremstyle{remark}
\newtheorem{remark}[theorem]{Remark}

\newcommand{\Z}{\mathbb Z}
\newcommand{\Xm}{X_m}
\newcommand{\epsb}{\varepsilon}

\title{Odd Moore Spaces Give a Projective-Free Answer\\
to a Hermiteness Question}
\author{Math discovery candidate full solution}
\date{August 17, 2026}

\begin{document}
\maketitle

\section{Status and source question}

Brudnyi and Sasane ask whether there is a Hausdorff space \(X\) with
\(H^n(X;\Z)\ne0\) for some \(n\ge5\) such that \(C(X)\) is Hermite
\cite{source}.  The question is printed on page 9 of the published-layout
arXiv PDF:

\begin{center}
\includegraphics[width=\textwidth]{figures/open_problem_crop.png}
\end{center}

The answer is yes.  In fact, a five-dimensional finite CW complex already
works, and its function algebra is projective free, which is stronger than
Hermite.  The same construction gives a sharp odd/even dichotomy.

\section{Proof intuition}

The source's sphere \(S^5\) fails because the exceptional group
\(\pi_4(U(2))\cong\Z/2\) produces a nontrivial stably trivial rank-two
bundle.  Replace the free top class of \(S^5\) by odd torsion: attach a
5-cell to \(S^4\) by an odd-degree map.  The universal coefficient theorem
moves the resulting \(\Z/m\) homology into \(H^5\), so the required
high-degree cohomology survives.

At the same time, the Puppe sequence says that vector bundles on this mapping
cone are controlled by multiplication by \(m\) on just two homotopy groups
of \(BU(r)\).  Odd multiplication is invertible on the troublesome
\(\Z/2\), while it is injective on the adjacent \(\Z\).  Exactness therefore
leaves no nontrivial bundle in any rank.

\section{The sharp Moore-space theorem}

For an integer \(m\ge2\), let
\[
 \Xm=M(\Z/m,4):=S^4\cup_{d_m} e^5,                            \tag{1}
\]
where \(d_m:S^4\to S^4\) has degree \(m\).

\begin{theorem}[Parity dichotomy]
\label{thm:main}
For every integer \(m\ge2\),
\[
                  H^5(\Xm;\Z)\cong\Z/m.
\]
Moreover:
\begin{enumerate}
\item if \(m\) is odd, every finite-rank complex vector bundle over \(\Xm\)
  is trivial, and \(C(\Xm)\) is projective free;
\item if \(m\) is even, \(\Xm\) supports a nontrivial stably trivial
  rank-two complex vector bundle, and \(C(\Xm)\) is not Hermite.
\end{enumerate}
\end{theorem}

\begin{proof}
\textbf{Cohomology.}
The reduced cellular chain complex of \(\Xm\) is concentrated in degrees
5 and 4:
\[
                 0\longrightarrow\Z
                   \xrightarrow{\ m\ }\Z\longrightarrow0.     \tag{2}
\]
Thus \(H_4(\Xm;\Z)\cong\Z/m\) and \(H_5(\Xm;\Z)=0\).  The integral
universal coefficient theorem gives
\[
 H^5(\Xm;\Z)
 \cong \operatorname{Ext}\bigl(H_4(\Xm;\Z),\Z\bigr)
 \cong \Z/m.                                                   \tag{3}
\]

\medskip
\noindent
\textbf{Classifying all complex vector bundles.}
The cofibration defining (1) continues as
\[
 S^4\xrightarrow{d_m}S^4\longrightarrow\Xm
 \longrightarrow S^5\xrightarrow{\Sigma d_m}S^5.             \tag{4}
\]
For \(Y=BU(r)\), its contravariant Puppe sequence contains the exact
segment of pointed sets
\[
 \pi_5(Y)\xrightarrow{\ m\ }\pi_5(Y)
 \longrightarrow[\Xm,Y]\longrightarrow
 \pi_4(Y)\xrightarrow{\ m\ }\pi_4(Y).                          \tag{5}
\]
Here precomposition by \(d_m\), or by its suspension, is multiplication by
\(m\) on the corresponding homotopy group.

The required low homotopy groups are
\[
\begin{array}{c|c|c}
r&\pi_5(BU(r))&\pi_4(BU(r))\\ \hline
1&0&0\\
2&\Z/2&\Z\\
r\ge3&0&\Z .
\end{array}                                                     \tag{6}
\]
Indeed, \(\pi_k(BU(r))=\pi_{k-1}(U(r))\);
\(\pi_4(U(2))=\pi_4(S^3)=\Z/2\), while
\(\pi_4(U(r))=0\) for \(r\ge3\), and
\(\pi_3(U(r))=\Z\) for \(r\ge2\) \cite{hatcherVB,mimura}.
Every \(BU(r)\) is simply connected, so for the connected complex \(\Xm\)
the pointed homotopy sets in (5) agree with the unpointed homotopy sets
that classify bundles.

Suppose first that \(m\) is odd.  In every row of (6), the first map in
(5) is surjective and the last map is injective.  Exactness forces
\([\Xm,BU(r)]\) to consist only of its basepoint.  Since rank-\(r\) complex
vector bundles over the finite CW complex \(\Xm\) are classified by
\([\Xm,BU(r)]\), every such bundle is trivial.

The space \(\Xm\) is connected and compact Hausdorff.  Proposition 2.1 of
the source identifies triviality of all its finite-rank complex vector
bundles with projective freeness of \(C(\Xm)\).  This proves part (1).

Now let \(m\) be even and set \(Y=BU(2)\).  In (5), multiplication by
\(m\) on \(\Z/2\) is zero, whereas multiplication by \(m\) on \(\Z\) is
injective.  Exactness shows that the middle pointed set has exactly two
elements:
\[
                    [\Xm,BU(2)]\cong\Z/2.                     \tag{7}
\]
The nonbasepoint class is a nontrivial rank-two bundle \(E\).

After stabilization to \(BU\), the corresponding Puppe sequence has
\(\pi_5(BU)=\pi_4(U)=0\) and \(\pi_4(BU)=\pi_3(U)=\Z\).  Hence
\([\Xm,BU]\) is a singleton, so the stabilized class of \(E\) is null.
Because \(\Xm\) is a finite CW complex, this nullhomotopy occurs at a finite
stage: for some \(k\),
\[
                       E\oplus\epsb^k\cong\epsb^{k+2}.         \tag{8}
\]
Thus \(E\) is nontrivial but stably trivial.  Proposition 2.3 of the source
then says that \(C(\Xm)\) is not Hermite.
\end{proof}

\begin{corollary}[Full affirmative answer]
For
\[
                     X=M(\Z/3,4)=S^4\cup_{d_3}e^5,
\]
one has \(H^5(X;\Z)\cong\Z/3\), and \(C(X)\) is projective free.  In
particular, \(C(X)\) is Hermite.
\end{corollary}

\begin{proof}
Apply the odd case of Theorem~\ref{thm:main} with \(m=3\).
\end{proof}

\section{Checks, scope, and novelty}

The pointed-set exactness in (5) is enough: injectivity of the last
multiplication makes every class restrict trivially, and surjectivity of the
first multiplication makes the preceding action have only the basepoint as
its image.  As an independent consistency check,
\[
 H^5(\Xm;\Z/2)
 \cong\operatorname{Ext}(\Z/m,\Z/2)
 =
 \begin{cases}
  0,&m\ \text{odd},\\
  \Z/2,&m\ \text{even},
 \end{cases}
\]
which exactly matches the exceptional \(\pi_4(U(2))\) term.

The result answers only the source's existence question; it does not
classify all Hausdorff spaces whose continuous-function algebras are
Hermite.  It does, however, classify the full Moore-space family (1) by
parity and proves the stronger projective-free conclusion in the affirmative
cases.

Bounded searches through August 17, 2026 used the exact question, exact
title, arXiv id, and combinations of Moore spaces, complex vector bundles,
Hermite algebras, and projective freeness.  The published 2023 paper still
states the question.  Its two citing works indexed by OpenAlex concern
bounded-variation and Dirichlet-series algebras, not this construction.  No
prior answer was located, so novelty is a candidate claim pending specialist
review.

\begin{thebibliography}{9}
\bibitem{source}
A. Brudnyi and A. Sasane,
\emph{Projective freeness and Hermiteness of complex function algebras},
Advances in Mathematics 434 (2023), 109322;
arXiv:2208.04901.

\bibitem{hatcherAT}
A. Hatcher,
\emph{Algebraic Topology},
Cambridge University Press, 2002.

\bibitem{hatcherVB}
A. Hatcher,
\emph{Vector Bundles and K-Theory},
Version 2.2, 2017.

\bibitem{mimura}
M. Mimura and H. Toda,
\emph{Topology of Lie Groups, I and II},
Translations of Mathematical Monographs 91,
American Mathematical Society, 1991.
\end{thebibliography}

\end{document}
